% CHECKLIST PROJET MASTER 2 - DT1 CAMEROUN
% Étudiante: Sorelle
% Encadrant: Dr. TCHAPET NJAFA Jean-Pierre

\documentclass[a4paper,11pt]{article}
\usepackage[utf8]{inputenc}
\usepackage[french]{babel}
\usepackage[margin=2cm]{geometry}
\usepackage{enumitem}
\usepackage{xcolor}
\usepackage{amssymb}
\usepackage{graphicx}
\usepackage{fancyhdr}
\usepackage{hyperref}

% Configuration
\hypersetup{
    colorlinks=true,
    linkcolor=blue,
    urlcolor=blue,
    pdftitle={Checklist Projet Master 2 - DT1 Cameroun},
    pdfauthor={Sorelle}
}

\pagestyle{fancy}
\fancyhf{}
\fancyhead[L]{\textbf{Projet Master 2 - DT1 Cameroun}}
\fancyhead[R]{Sorelle - 2025/2026}
\fancyfoot[C]{\thepage}

% Commandes personnalisées
\newcommand{\checkbox}{\large$\square$}
\newcommand{\checkboxdone}{\large$\boxtimes$}
\newcommand{\priorite}[1]{\textcolor{red}{\textbf{[#1]}}}
\newcommand{\duree}[1]{\textcolor{blue}{\textit{(#1)}}}

\title{\textbf{\Huge CHECKLIST PROJET}\\[0.5cm]
       \Large Détection précoce du Diabète Type 1\\
       au Cameroun par Machine Learning\\[0.3cm]
       \large Master 2 Biophysique}
\author{Étudiante: \textbf{Sorelle}\\
        Encadrant: Dr. TCHAPET NJAFA Jean-Pierre\\
        Conseiller: Prof. NANA ENGO Serge Guy}
\date{Novembre 2025 - Mai 2026}

\begin{document}

\maketitle
\thispagestyle{empty}

\vspace{1cm}

\begin{center}
\fbox{\parbox{0.9\textwidth}{
\centering
\textbf{INSTRUCTIONS D'UTILISATION}\\[0.3cm]
\begin{itemize}[leftmargin=*]
    \item Imprimez ce document et cochez les cases au fur et à mesure
    \item \priorite{P1} = Priorité critique | \priorite{P2} = Important | \priorite{P3} = Optionnel
    \item \duree{Xj/Xs} = Durée estimée (jours/semaines)
    \item Consultez le diagramme de Gantt pour la planification détaillée
\end{itemize}
}}
\end{center}

\newpage

\section*{PHASE 1: INSTALLATION ET PRISE EN MAIN \duree{1-2 semaines}}

\subsection*{1.1 Installation de l'environnement \priorite{P1}}

\begin{itemize}[label=\checkboxdone]
    \item Installer Python 3.9+ (Anaconda recommandé)
    \item Créer l'environnement virtuel: \texttt{conda create -n dt1\_cameroun python=3.9}
    \item Activer l'environnement: \texttt{conda activate dt1\_cameroun}
    \item Installer les dépendances: \texttt{pip install -r requirements\_full.txt}
    \item Tester l'installation: \texttt{python -c "import pandas; import sklearn; print('OK')"}
    \item Installer Jupyter: \texttt{pip install jupyter notebook}
    \item Lancer Jupyter: \texttt{jupyter notebook}
\end{itemize}

\subsection*{1.2 Découverte du projet \priorite{P1}}

\begin{itemize}[label=\checkboxdone]
    \item Lire \texttt{README.md}
    \item Lire \texttt{GUIDE\_EXECUTION.md}
    \item Lire \texttt{demarche\_methodologique\_sorelle.md}
    \item Explorer la structure des dossiers
    \item Vérifier que tous les notebooks s'ouvrent correctement
\end{itemize}

\subsection*{1.3 Gestion bibliographique \priorite{P1}}

\begin{itemize}[label=\checkboxdone]
    \item Installer Zotero: \url{https://www.zotero.org/download/}
    \item Créer un compte Zotero
    \item Importer les 36 PDFs dans Zotero (glisser-déposer)
    \item Vérifier les métadonnées (titre, auteurs, année, journal)
    \item Créer des collections par thème: DT1\_Afrique, Biomarqueurs, ML\_Algo, etc.
    \item Installer le plugin Zotero pour LaTeX
    \item Exporter la bibliothèque: Fichier → Exporter → Format BibTeX
    \item Sauvegarder dans: \texttt{6\_MEMOIRE\_LATEX/bibliographie.bib}
\end{itemize}

\newpage

\section*{PHASE 2: APPRENTISSAGE MACHINE LEARNING \duree{7 semaines}}

\subsection*{2.1 Semaine 1: Introduction Python + ML \priorite{P1} \duree{1 semaine}}

\begin{itemize}[label=\checkboxdone]
    \item Ouvrir le notebook: \texttt{Semaine01\_Introduction\_Python\_ML.ipynb}
    \item \textbf{Section 1:} Comprendre NumPy (vecteurs, matrices)
    \item \textbf{Section 2:} Maîtriser Pandas (DataFrame, manipulation)
    \item \textbf{Section 3:} Visualisation avec Matplotlib
    \item \textbf{Section 4:} Premier modèle (Decision Tree)
    \item \textbf{Exercices:} Compléter les 5 exercices pratiques
    \item \textbf{Auto-évaluation:} Vérifier les réponses
    \item \textbf{Validation:} Exécuter toutes les cellules sans erreur
\end{itemize}

\subsection*{2.2 Semaine 2: Exploration des données \priorite{P1} \duree{1 semaine}}

\begin{itemize}[label=\checkboxdone]
    \item Ouvrir: \texttt{Semaine02\_Exploration\_Donnees.ipynb}
    \item Générer le dataset: \texttt{python 2\_DONNEES/generate\_synthetic\_data.py}
    \item Vérifier: \texttt{2\_DONNEES/raw/dataset\_dt1\_cameroun\_synthetic.csv}
    \item Charger et inspecter les données (1000 patients)
    \item Statistiques descriptives (moyennes, écarts-types)
    \item Tests statistiques (Shapiro-Wilk, Mann-Whitney)
    \item Matrice de corrélation
    \item Visualisations (distributions, boxplots, heatmap)
    \item Comprendre les biomarqueurs: ANP32A-IT1, ESCO2, NBPF1
    \item \textbf{Exercices:} Analyser le déséquilibre des classes
\end{itemize}

\subsection*{2.3 Semaine 3: Modèles simples \priorite{P1} \duree{1 semaine}}

\begin{itemize}[label=\checkboxdone]
    \item Ouvrir: \texttt{Semaine03\_Modelisation\_Simple.ipynb}
    \item \textbf{Modèle 1:} Logistic Regression
    \item \textbf{Modèle 2:} Decision Tree
    \item \textbf{Modèle 3:} Random Forest (100 arbres)
    \item \textbf{Modèle 4:} Support Vector Machine (SVM)
    \item Calculer les métriques: Accuracy, Precision, Recall, F1, AUC
    \item Comparer les 4 modèles (tableau + graphiques)
    \item Tracer les courbes ROC
    \item Identifier le meilleur modèle selon le Recall
    \item \textbf{Comprendre:} Pourquoi prioriser le Recall en médecine?
\end{itemize}

\subsection*{2.4 Semaines 4-5: Modèles avancés \priorite{P1} \duree{2 semaines}}

\begin{itemize}[label=\checkbox]
    \item Ouvrir: \texttt{Semaine04-05\_Modeles\_Avances.ipynb}
    \item \textbf{XGBoost:} Version de base
    \item \textbf{XGBoost:} Avec \texttt{scale\_pos\_weight} (gérer déséquilibre)
    \item \textbf{LightGBM:} Avec \texttt{class\_weight='balanced'}
    \item \textbf{SMOTE:} Génération d'échantillons synthétiques DT1
    \item Random Forest + SMOTE
    \item Comparer les performances (Recall prioritaire)
    \item Analyser les Feature Importances
    \item Comprendre le trade-off Precision/Recall
    \item Identifier les 3 biomarqueurs les plus importants
\end{itemize}

\subsection*{2.5 Semaine 6: Optimisation \priorite{P1} \duree{1 semaine}}

\begin{itemize}[label=\checkbox]
    \item Ouvrir: \texttt{Semaine06\_Optimisation\_Validation.ipynb}
    \item Validation croisée (Stratified K-Fold, k=5)
    \item GridSearchCV sur Random Forest (exhaustif)
    \item RandomizedSearchCV sur XGBoost (50 combinaisons)
    \item Comparer les deux approches
    \item Sélectionner le modèle final (meilleur Recall)
    \item Sauvegarder le modèle: \texttt{4\_MODELES/models/final\_model.pkl}
    \item Sauvegarder les features: \texttt{4\_MODELES/models/features.pkl}
    \item Tester le chargement du modèle
    \item Créer le rapport final du modèle
\end{itemize}

\subsection*{2.6 Semaine 7: Interprétabilité \priorite{P1} \duree{1 semaine}}

\begin{itemize}[label=\checkbox]
    \item Ouvrir: \texttt{Semaine07\_Interpretation\_Resultats.ipynb}
    \item Charger le modèle final sauvegardé
    \item Calculer SHAP values (TreeExplainer)
    \item Générer Summary Plot (importance globale)
    \item Générer Dependence Plots (effet de chaque feature)
    \item Générer Force Plots (explications individuelles)
    \item Analyser les erreurs de prédiction (Faux Positifs, Faux Négatifs)
    \item Créer les figures pour le mémoire (haute résolution)
    \item Sauvegarder dans: \texttt{6\_MEMOIRE\_LATEX/figures/}
    \item Interpréter cliniquement les résultats
\end{itemize}

\newpage

\section*{PHASE 3: LECTURES BIBLIOGRAPHIQUES \duree{Parallèle aux notebooks}}

\subsection*{3.1 Articles prioritaires DT1 Afrique \priorite{P1} \duree{2 semaines}}

\begin{itemize}[label=\checkbox]
    \item \textbf{[CRITIQUE]} Katte 2023 - Phénotype DT1 en Afrique subsaharienne
    \item Oram 2025 - Impact ancestry sur GRS (medRxiv)
    \item Kengne 2013 - Diabète et obésité en Afrique (Heart/BMJ)
    \item Prendre des notes dans Zotero (annotations PDF)
    \item Rédiger fiche de lecture pour chaque article (1 page)
\end{itemize}

\subsection*{3.2 Articles biomarqueurs \priorite{P1} \duree{1 semaine}}

\begin{itemize}[label=\checkbox]
    \item \textbf{[CRITIQUE]} Endocrine Journal 2023 - ANP32A-IT1, ESCO2, NBPF1
    \item Comprendre les 3 biomarqueurs et leur AUC > 0.8
    \item Noter les valeurs seuils et performances diagnostiques
\end{itemize}

\subsection*{3.3 Articles ML médical \priorite{P2} \duree{2 semaines}}

\begin{itemize}[label=\checkbox]
    \item Contreras 2018 - IA pour le diabète (JMIR)
    \item Sendak 2020 - ML en contexte clinique réel (arXiv)
    \item Comprendre les défis du déploiement clinique
\end{itemize}

\subsection*{3.4 Articles techniques ML \priorite{P2} \duree{2 semaines}}

\begin{itemize}[label=\checkbox]
    \item Chen 2016 - XGBoost (algorithme principal)
    \item Ke 2017 - LightGBM
    \item Breiman 2001 - Random Forests
    \item Lundberg 2017 - SHAP (interprétabilité)
    \item Chawla 2002 - SMOTE (déséquilibre)
\end{itemize}

\subsection*{3.5 Livres de référence \priorite{P3} \duree{En continu}}

\begin{itemize}[label=\checkbox]
    \item Hastie 2009 - Elements of Statistical Learning (chapitres 9, 10, 15)
    \item James 2021 - Intro Statistical Learning (chapitres 4, 5, 8)
    \item Parcourir les chapitres pertinents (pas lecture intégrale)
\end{itemize}

\newpage

\section*{PHASE 4: APPLICATION STREAMLIT \duree{3 semaines}}

\subsection*{4.1 Tests de l'application \priorite{P1} \duree{1 semaine}}

\begin{itemize}[label=\checkbox]
    \item Lancer l'app: \texttt{streamlit run 5\_APPLICATION\_DEMO/app\_streamlit.py}
    \item Tester la page d'accueil
    \item Tester la prédiction manuelle (saisie des 7 features)
    \item Tester l'upload CSV (fichier d'exemple)
    \item Vérifier les prédictions et probabilités
    \item Identifier les bugs éventuels
\end{itemize}

\subsection*{4.2 Ajout de fonctionnalités \priorite{P2} \duree{1 semaine}}

\begin{itemize}[label=\checkbox]
    \item Intégrer le modèle final optimisé (remplacer modèle de démo)
    \item Ajouter graphiques SHAP pour explications
    \item Ajouter statistiques sur les prédictions batch
    \item Améliorer les messages d'erreur
    \item Ajouter validation des inputs
\end{itemize}

\subsection*{4.3 Finalisation interface \priorite{P2} \duree{1 semaine}}

\begin{itemize}[label=\checkbox]
    \item Améliorer le design (couleurs, mise en page)
    \item Ajouter logo/bannière
    \item Ajouter page "À propos" détaillée
    \item Documenter l'utilisation (README dans dossier app)
    \item Tester sur différents navigateurs
    \item Créer capture d'écran pour le mémoire
\end{itemize}

\newpage

\section*{PHASE 5: RÉDACTION DU MÉMOIRE \duree{12 semaines}}

\subsection*{5.1 Introduction \priorite{P1} \duree{2 semaines}}

\begin{itemize}[label=\checkbox]
    \item Lire \texttt{01\_introduction.tex} (déjà pré-remplie)
    \item Compléter le contexte épidémiologique (chiffres IDF)
    \item Développer la problématique du dataset shift
    \item Justifier le choix des 3 biomarqueurs
    \item Formuler clairement les objectifs
    \item Annoncer le plan du mémoire
    \item Relire et corriger (orthographe, syntaxe)
    \item Faire relire par l'encadrant
\end{itemize}

\subsection*{5.2 État de l'art \priorite{P1} \duree{3 semaines}}

\begin{itemize}[label=\checkbox]
    \item \textbf{Section 2.1:} DT1 en Afrique (Katte, Kengne, etc.)
    \item \textbf{Section 2.2:} Biomarqueurs génétiques (Endocrine J., Oram)
    \item \textbf{Section 2.3:} ML en santé (Contreras, Sendak)
    \item \textbf{Section 2.4:} Algorithmes ML (XGBoost, RF, SMOTE)
    \item \textbf{Section 2.5:} Interprétabilité (SHAP, LIME)
    \item \textbf{Section 2.6:} Défis et limites
    \item Citer 40-50 références
    \item Vérifier cohérence des citations LaTeX
    \item Créer tableau récapitulatif des études
\end{itemize}

\subsection*{5.3 Méthodologie \priorite{P1} \duree{3 semaines}}

\begin{itemize}[label=\checkbox]
    \item \textbf{Section 3.1:} Description du dataset (1000 patients synthétiques)
    \item \textbf{Section 3.2:} Pré-traitement et exploration
    \item \textbf{Section 3.3:} Algorithmes testés (4 de base + 3 avancés)
    \item \textbf{Section 3.4:} Gestion du déséquilibre (SMOTE, class\_weight)
    \item \textbf{Section 3.5:} Optimisation (GridSearch, RandomSearch)
    \item \textbf{Section 3.6:} Validation croisée (Stratified K-Fold)
    \item \textbf{Section 3.7:} Métriques d'évaluation (Recall prioritaire)
    \item \textbf{Section 3.8:} Interprétabilité (SHAP)
    \item Créer schéma méthodologique (workflow)
    \item Créer tableau des hyperparamètres
\end{itemize}

\subsection*{5.4 Résultats \priorite{P1} \duree{2 semaines}}

\begin{itemize}[label=\checkbox]
    \item \textbf{Générer TOUTES les figures} (depuis notebooks)
    \item Figure 1: Distribution des biomarqueurs (boxplots)
    \item Figure 2: Matrice de corrélation
    \item Figure 3: Comparaison des modèles (barres)
    \item Figure 4: Courbes ROC comparatives
    \item Figure 5: Feature Importances (XGBoost final)
    \item Figure 6: SHAP Summary Plot
    \item Figure 7: SHAP Dependence Plots (top 3 features)
    \item Tableau 1: Caractéristiques du dataset
    \item Tableau 2: Performances des modèles (Accuracy, Precision, Recall, F1, AUC)
    \item Tableau 3: Hyperparamètres optimaux du modèle final
    \item Tableau 4: Matrice de confusion
    \item \textbf{Rédiger:} Présenter les résultats sans interprétation
    \item Numéroter et légender toutes les figures/tableaux
\end{itemize}

\subsection*{5.5 Discussion \priorite{P1} \duree{2 semaines}}

\begin{itemize}[label=\checkbox]
    \item Interpréter les performances (Recall 85-90\% attendu)
    \item Discuter l'importance des biomarqueurs (ESCO2 > ANP32A-IT1 > NBPF1?)
    \item Comparer avec la littérature (Endocrine J. 2023)
    \item Discuter le déséquilibre des classes (impact SMOTE)
    \item Interpréter SHAP (contribution de chaque feature)
    \item Discuter les erreurs (FP vs FN, analyse clinique)
    \item Discuter les limites (données synthétiques, dataset shift)
    \item Discuter la transférabilité au contexte camerounais
    \item Discuter l'interprétabilité pour le clinicien
    \item Proposer des perspectives d'amélioration
\end{itemize}

\subsection*{5.6 Conclusion + Abstract \priorite{P1} \duree{1 semaine}}

\begin{itemize}[label=\checkbox]
    \item Rappeler le contexte et les objectifs
    \item Synthétiser les résultats principaux
    \item Confirmer/infirmer les hypothèses
    \item Apports du travail (scientifiques, méthodologiques)
    \item Limites reconnues
    \item Perspectives (validation sur données réelles, déploiement)
    \item Rédiger l'Abstract (200-250 mots, français + anglais)
    \item Mots-clés (5-7): Diabète Type 1, Machine Learning, XGBoost, Biomarqueurs, etc.
\end{itemize}

\subsection*{5.7 Annexes et compléments \priorite{P2} \duree{1 semaine}}

\begin{itemize}[label=\checkbox]
    \item Annexe A: Code source principal (extraits commentés)
    \item Annexe B: Tableaux de résultats détaillés
    \item Annexe C: Figures complémentaires
    \item Liste des abréviations
    \item Liste des figures
    \item Liste des tableaux
    \item Table des matières (automatique LaTeX)
\end{itemize}

\newpage

\section*{PHASE 6: RÉVISION ET FINALISATION \duree{4 semaines}}

\subsection*{6.1 Relecture et corrections \priorite{P1} \duree{2 semaines}}

\begin{itemize}[label=\checkbox]
    \item Relecture complète (orthographe, grammaire, syntaxe)
    \item Vérifier cohérence des citations (bibliographie)
    \item Vérifier numérotation figures/tableaux
    \item Vérifier les références croisées LaTeX
    \item Harmoniser la terminologie
    \item Vérifier les formules mathématiques
    \item Soumettre à l'encadrant pour première relecture
    \item Intégrer les corrections de l'encadrant
    \item Soumettre au conseiller pour relecture
    \item Intégrer les corrections du conseiller
\end{itemize}

\subsection*{6.2 Mise en page finale \priorite{P1} \duree{1 semaine}}

\begin{itemize}[label=\checkbox]
    \item Compiler LaTeX: \texttt{pdflatex main.tex} (3 fois)
    \item Compiler BibTeX: \texttt{bibtex main}
    \item Vérifier PDF final (pagination, marges, qualité figures)
    \item Vérifier que toutes les figures sont en haute résolution
    \item Ajuster les sauts de page si nécessaire
    \item Vérifier les en-têtes et pieds de page
    \item Générer le PDF final: \texttt{memoire\_sorelle\_dt1\_2026.pdf}
    \item Vérifier la taille (<60 pages hors annexes)
\end{itemize}

\subsection*{6.3 Impression et reliure \priorite{P1} \duree{3 jours}}

\begin{itemize}[label=\checkbox]
    \item Imprimer 1 exemplaire test (vérifier qualité)
    \item Corriger si nécessaire
    \item Imprimer 4 exemplaires définitifs (jury + encadrant + vous)
    \item Relier les exemplaires (reliure spirale ou thermique)
    \item Graver 2 CD/USB avec: PDF + code source + données + notebooks
    \item Préparer pochette avec nom, titre, année
\end{itemize}

\subsection*{6.4 Dépôt du mémoire \priorite{P1} \duree{1 jour}}

\begin{itemize}[label=\checkbox]
    \item Vérifier la date limite de dépôt (secrétariat)
    \item Déposer les exemplaires au secrétariat
    \item Déposer la version numérique (email ou plateforme)
    \item Obtenir le récépissé de dépôt
    \item Envoyer copie PDF à l'encadrant et conseiller
\end{itemize}

\newpage

\section*{PHASE 7: SOUTENANCE \duree{2 semaines}}

\subsection*{7.1 Préparation des slides \priorite{P1} \duree{1 semaine}}

\begin{itemize}[label=\checkbox]
    \item Créer présentation PowerPoint/Beamer (15-20 slides)
    \item \textbf{Slide 1:} Page de titre (nom, titre, date, jury)
    \item \textbf{Slide 2:} Plan de la présentation
    \item \textbf{Slides 3-4:} Contexte et problématique
    \item \textbf{Slide 5:} Objectifs
    \item \textbf{Slides 6-8:} Méthodologie (schéma workflow)
    \item \textbf{Slides 9-12:} Résultats (figures principales)
    \item \textbf{Slide 13:} Discussion et interprétation
    \item \textbf{Slide 14:} Limites et perspectives
    \item \textbf{Slide 15:} Conclusion
    \item \textbf{Slide 16:} Remerciements
    \item Vérifier lisibilité (police >= 20pt, pas trop de texte)
    \item Ajouter numéros de slides
    \item Faire valider par l'encadrant
\end{itemize}

\subsection*{7.2 Répétitions \priorite{P1} \duree{1 semaine}}

\begin{itemize}[label=\checkbox]
    \item Répétition 1: Seul(e), chronométrer (20 min)
    \item Répétition 2: Devant l'encadrant (feedback)
    \item Répétition 3: Devant collègues/famille (questions)
    \item Préparer réponses aux questions prévisibles:
    \item \quad - Pourquoi ces biomarqueurs?
    \item \quad - Pourquoi XGBoost plutôt que RF?
    \item \quad - Comment gérer le dataset shift?
    \item \quad - Validation sur données réelles?
    \item \quad - Déploiement clinique?
    \item Répétition 4: Conditions réelles (amphi, projecteur)
    \item Préparer tenue professionnelle
\end{itemize}

\subsection*{7.3 Jour de la soutenance \priorite{P1}}

\begin{itemize}[label=\checkbox]
    \item Arriver 30 min en avance
    \item Tester le projecteur et l'ordinateur
    \item Avoir une clé USB de secours (slides + PDF mémoire)
    \item Saluer le jury
    \item \textbf{PRÉSENTATION:} 20 minutes
    \item \textbf{QUESTIONS:} 20-30 minutes
    \item Rester calme et confiant(e)
    \item Remercier le jury
    \item \textbf{FÉLICITATIONS! Vous êtes désormais titulaire d'un Master 2!}
\end{itemize}

\newpage

\section*{ANNEXES}

\subsection*{Contacts importants}

\begin{itemize}
    \item \textbf{Encadrant:} Dr. TCHAPET NJAFA Jean-Pierre - [email à compléter]
    \item \textbf{Conseiller:} Prof. NANA ENGO Serge Guy - [email à compléter]
    \item \textbf{Secrétariat Master:} [contact à compléter]
    \item \textbf{Bibliothèque universitaire:} [contact à compléter]
\end{itemize}

\subsection*{Ressources utiles}

\begin{itemize}
    \item \textbf{Documentation Python:} \url{https://docs.python.org/3/}
    \item \textbf{Scikit-learn:} \url{https://scikit-learn.org/}
    \item \textbf{XGBoost:} \url{https://xgboost.readthedocs.io/}
    \item \textbf{SHAP:} \url{https://shap.readthedocs.io/}
    \item \textbf{Streamlit:} \url{https://docs.streamlit.io/}
    \item \textbf{LaTeX:} \url{https://www.overleaf.com/learn}
    \item \textbf{Zotero:} \url{https://www.zotero.org/support/}
\end{itemize}

\subsection*{Statistiques du projet}

\begin{center}
\begin{tabular}{|l|r|}
\hline
\textbf{Élément} & \textbf{Quantité} \\
\hline
Durée totale & 6 mois (Nov 2025 - Mai 2026) \\
Notebooks pédagogiques & 6 \\
Articles scientifiques & 36 PDFs (79 MB) \\
Références bibliographiques & 64+ \\
Lignes de code Python & $\sim$2700 \\
Dataset (patients) & 1000 \\
Features (variables) & 7 \\
Algorithmes testés & 7 \\
Pages mémoire (max) & 60 \\
\hline
\end{tabular}
\end{center}

\vspace{2cm}

\begin{center}
\fbox{\parbox{0.9\textwidth}{
\centering
\textbf{BON COURAGE!}\\[0.3cm]
Ce projet est ambitieux mais parfaitement réalisable en suivant cette checklist.\\
Travaillez régulièrement, n'hésitez pas à demander de l'aide,\\
et célébrez chaque étape franchie!\\[0.5cm]
\textit{``Le succès est la somme de petits efforts répétés jour après jour.''}\\
-- Robert Collier
}}
\end{center}

\end{document}
